
En
appliquant 3.7, on trouve :
\[
\ch(g_*(f_*(x)))=
g_{\mathrm{gr}}\!\left(
 f_{\mathrm{gr}}\bigl(\ch(x)\cdot\Todd(\check T_f)\cdot
 f^{\mathrm{gr}}(\Todd(\check T_g))\bigr)
\right).
\]
La fonctorialité de \(c\) et de \(\Todd\) entraîne :
\[
 f^{\mathrm{gr}}(\Todd(\check T_g))=\Todd(f^*(\check T_g)).
\]
D'après 2.7, on a :
\[
 \check T_{g\circ f}=\check T_f+f^*(\check T_g).
\]
La multiplicativité du caractère de Todd entraîne donc :
\[
 \Todd(\check T_f)\cdot f^{\mathrm{gr}}(\Todd(\check T_g))
 =
 \Todd(\check T_{g\circ f}).
\]
On obtient donc finalement :
\[
 \ch(g_*(f_*(x)))=
 g_{\mathrm{gr}}\bigl(f_{\mathrm{gr}}\bigl(\ch(x)\cdot\Todd(\check T_{g\circ f})\bigr)\bigr),
\]
ce qui est le résultat annoncé.

\textbf{3.9.} Si on applique 3.8 à la factorisation de \(f\) donnée en 3.1, on voit qu'il suffit, pour démontrer le théorème de Riemann--Roch, de le démontrer dans le cas d'une immersion régulière de codimension constante, et dans le cas du morphisme structural : $\mathbf{P}^r_Y\longrightarrow Y$.

En outre, on peut supposer que l'immersion \(i\) satisfait la condition de VII 4.1. Soit en effet \(s\) la section de \(\mathbf{P}^2_{\mathbf{P}^r_Y}\) introduite en VII 4.2. Alors \(s\circ i\) est une immersion régulière vérifiant la condition de VII 4.1 (VII 4.2) et d'après 3.8, il suffit de montrer 3.6 pour \(s\circ i\), \(\mathbf{P}^2_{\mathbf{P}^r_Y}\longrightarrow\mathbf{P}^r_Y\) et \(\mathbf{P}^r_Y\longrightarrow Y\).

On va donc prouver 3.6 en prouvant les deux cas particuliers suivants :
\begin{enumerate}[label=\roman*)]
\item \(f\) est une immersion régulière de codimension constante, satisfaisant la condition de VII 4.1;
